% STALL REPORT: G2 (literal fixed-anchor anti-concentration) is proved below,
% but it has the wrong quantifier order for the scale-dependent anchored star.
% G1 is not proved; literal G2 is proved but does not imply AMTD; G3 is not
% reached; G4 is the principal output.  B1 and B2 as stated in
% CODEX_SPEC_oracleB.md confuse
% coefficient/Mellin zeros with evaluation roots.  The exact surviving
% arithmetic input is the marked two-characteristic crystalline Mellin
% dispersion statement (eq:oracleB-mh2), which remains open.

\subsection{Arithmetic oracle B: a fixed-anchor theorem and the sharp
horizontal obstruction}
\label{ssec:oracleB-result}

Put
\[
 \mathcal Z_p=\{0\leq j<p:b_j\equiv0\pmod p\},
 \qquad H_p(t)=\sum_{j=0}^{p-1}b_jt^j.
\]
The requested AMTD estimate is not proved here.  There is nevertheless an
unconditional answer to the literal fixed-target question in goal~G2.  It is
stronger than density zero, but it also exposes the decisive quantifier
reversal: a target fixed before the prime varies is elementary, whereas the
no-go target is chosen after the dyadic scale.
Thus the status of the ranked goals is: G1 is open; literal G2 is proved
below but is insufficient for the moving target; G3 is not achieved because
the required marked crystalline coordinate has not been constructed; and G4,
the sharp obstruction with the exact missing input, is achieved.

\begin{proposition}[Fixed Ap\'ery anchors are finite]
\label{prop:oracleB-fixed-anchor}
For $c\in\mathbb Z$, let $\bar c_p\in\{0,\ldots,p-1\}$ be its residue
modulo~$p$, and set
\[
 c^\flat=\begin{cases}
 c,&c\geq0,\\
 -c-1,&c<0.
 \end{cases}
\]
Then
\begin{equation}
 \{p\geq5:\bar c_p\in\mathcal Z_p\}
 \subseteq \{p\geq5:p\mid b_{c^\flat}\}.
 \label{eq:oracleB-fixed-inclusion}
\end{equation}
Consequently, for every fixed~$c$,
\begin{equation}
 \#\{p\leq X:\bar c_p\in\mathcal Z_p\}
 \leq\omega(b_{c^\flat})=O_c(1)=o(\pi(X)).
 \label{eq:oracleB-fixed-density}
\end{equation}
For $p>|c|$, the inclusion in~\eqref{eq:oracleB-fixed-inclusion} is an
equivalence.
\end{proposition}

\begin{proof}
First suppose that $c\geq0$, and write $c=\sum_i c_ip^i$ in base~$p$.
The lowest digit is $c_0=\bar c_p$.  If
$\bar c_p\in\mathcal Z_p$, then $p\mid b_{c_0}$, and Gessel's
Lucas congruence (Theorem~\ref{thm:lucas}) gives
\[
 b_c\equiv\prod_i b_{c_i}\equiv0\pmod p.
\]

Now let $c=-d<0$ and write $d=qp+s$ with $0\leq s<p$.  If $s=0$,
then $\bar c_p=0$ and there is no hit because $b_0=1$.  If $s>0$,
then $\bar c_p=p-s$.  Reflection (Remark~\ref{rem:orbit}, or
Proposition~\ref{prop:reflection}) gives
\[
 b_{p-s}\equiv b_{s-1}\pmod p.
\]
Thus a hit implies $p\mid b_{s-1}$.  The lowest base-$p$ digit of
$d-1$ is $s-1$, so another application of Lucas gives
$p\mid b_{d-1}=b_{c^\flat}$.  This proves the inclusion.

The binomial formula for $b_n$ has positive integral summands and its
$k=0$ summand is~$1$; hence $b_n$ is a fixed nonzero integer for fixed
$n$.  It has only finitely many prime divisors, proving
\eqref{eq:oracleB-fixed-density}.  Finally, if $p>|c|$, then for
$c\geq0$ one has $\bar c_p=c$, while for $c=-d<0$ reflection gives
$b_{\bar c_p}=b_{p-d}\equiv b_{d-1}$; both implications are then
reversible.
\end{proof}

\begin{remark}[Why this does not kill the anchored star]
\label{rem:oracleB-quantifiers}
Proposition~\ref{prop:oracleB-fixed-anchor} has the quantifier order
\begin{equation}
 \forall c\in\mathbb Z\;\exists C_c\;\forall X:\qquad
 \#\{p\leq X:\bar c_p\in\mathcal Z_p\}\leq C_c.
 \label{eq:oracleB-fixed-quantifiers}
\end{equation}
In Section~\ref{ssec:nogo}, however, the anchor is
$m_0=m_0(N)\in I_N$ and is allowed to move with~$N$.  To exclude just
one full star in a dyadic prime block, one would already need a uniform
statement such as
\begin{equation}
 \sup_{m\in I_N}
 \#\{P<p\leq2P:m\bmod p\in\mathcal Z_p\}
   =o(P/\log P),
 \label{eq:oracleB-moving-anchor}
\end{equation}
in the relevant range $P>\sqrt N$.  AMTD asks for a centered second
moment, which is stronger still.

Lucas merely turns the expression inside the supremum into a count of
large prime divisors of the moving integer~$b_m$.  The elementary height
bound $\log b_m=O(m)$ gives at best $O(N/\log N)$ prime divisors of size
$\asymp N$ when $m\asymp N$, exactly the size of a full prime block.
Thus this Lucas/height argument supplies no way to move the constants
$C_c$ outside the supremum.  This is the same distinction exhibited
by the anchored-star/no-go construction in
Section~\ref{sec:second-moment}.
\end{remark}

\subsubsection{The two zero loci are in dual spaces}

The proposed B1 and B2 routes make an object-level identification that is
false.  The following elementary calculation makes the distinction exact.

\begin{proposition}[Coefficient zeros are not Hasse-polynomial roots]
\label{prop:oracleB-two-zero-loci}
For $1\leq j\leq p-2$, the following identity holds in~$\mathbb F_p$:
\begin{equation}
 b_j=-\sum_{t\in\mathbb F_p^\times}H_p(t)t^{-j}.
 \label{eq:oracleB-mellin}
\end{equation}
Consequently,
\[
 \mathcal Z_p=\{j:[t^j]H_p(t)=0\}
\]
is a coefficient, or finite-Mellin, zero-set.  It is not the evaluation
root locus of $H_p$ or of the square root in
Remark~\ref{rem:squareness}.
\end{proposition}

\begin{proof}
After inserting $H_p(t)=\sum_{k=0}^{p-1}b_kt^k$, multiplicative
orthogonality gives
\[
 \sum_{t\in\mathbb F_p^\times}H_p(t)t^{-j}
 =\sum_{k=0}^{p-1}b_k\sum_{t\in\mathbb F_p^\times}t^{k-j}
 =-b_j,
\]
because $1\leq j\leq p-2$ and the inner sum vanishes unless $k=j$.

For a concrete separation, exact reduction at $p=7$ gives
\begin{align*}
 H_7(t)
 &=1+5t+3t^2+3t^3+3t^4+5t^5+t^6\\
 &=(t-1)^2(t^2+1)^2\pmod7.
\end{align*}
No coefficient is zero, so $\mathcal Z_7=\varnothing$, whereas $t=1$
is an evaluation root of both $H_7$ and its square root.
\end{proof}

In particular, if $\rho_p=\min\mathcal Z_p$, then $\rho_p$ need not be
a root of~$B_p(t)$, and no theorem identifies the two notions; they can
coincide accidentally at individual primes.  A Weil bound for
$V(B_p)$ therefore does not bound
\[
 F_p(a)=\sum_{j\in\mathcal Z_p}e_p(aj).
\]
Nor is $\rho_p$ equivalent to the whole zero-set: it is undefined when
$\mathcal Z_p$ is empty, and for example
\[
 \mathcal Z_{181}=\{19,47,133,161\}.
\]
For a center singleton the first zero trivially determines the set; only on
the doublet subfamily $Z(p)=2$ does it determine a noncentral reflected pair
and its gap.  Once there are at least two noncentral pairs, it does not
determine the whole zero-set.

\subsubsection{Why the square root is not a bounded-complexity family}

Let $\Delta(t)=t^2-34t+1$.  The factorization stated in
Remark~\ref{rem:squareness} and proved, at the level of factorization, by
Caruso--F\"urnsinn--Vargas-Montoya--Zudilin~\cite{CFVZ2026} is
\begin{equation}
 H_p(t)=\Delta(t)^{\varepsilon_p}B_p(t)^2,
 \qquad
 \deg B_p=\frac{p-1-2\varepsilon_p}{2}.
 \label{eq:oracleB-square}
\end{equation}
Thus $\deg B_p=\Theta(p)$.  This yields a second, independent obstruction
to B2 even if one temporarily replaces coefficient zeros by evaluation
roots.

\begin{proposition}[Fixed family versus moving Hasse section]
\label{prop:oracleB-complexity}
The polynomials $B_p$ cannot be the reductions of one fixed bounded-degree
divisor on the $t$-line.  If $B_p$ is squarefree, then
$V(B_p)_{\overline{\mathbb F}_p}$ has $\deg B_p=\Theta(p)$ geometric
points.  The equation $B_p(t)=0$ is zero-dimensional and has no genus;
if it is replaced by the hyperelliptic curve $y^2=B_p(t)$, then
\[
 g_p=\left\lfloor\frac{\deg B_p-1}{2}\right\rfloor=\Theta(p).
\]
Hence the standard degree/genus form of the Weil--Bombieri estimate supplies
no constant uniform in~$p$ for either version.
\end{proposition}

\begin{proof}
A fixed relative divisor of degree~$D$ has fibers of degree~$D$, whereas
\eqref{eq:oracleB-square} has unbounded degree.  The remaining assertions
are the point count of a squarefree zero-dimensional divisor and the
standard genus formula for a squarefree hyperelliptic equation.
\end{proof}

There is no contradiction between this proposition and the existence of a
fixed Picard--Fuchs family.  A fixed family can produce, separately in each
characteristic~$p$, a new Frobenius/Hasse section whose weight and degree
grow with~$p$.  ``The geometric family is fixed'' does not imply ``the
derived Hasse polynomial has bounded complexity.''  This is precisely the
moving-polynomial versus fixed-family quantifier distinction.

There are two source cautions.  First, Theorem~1.2 of~\cite{CFVZ2026}
states the square/quadratic-times-square factorization.  Its proof uses the
Franel relation
\begin{equation}
 f_\alpha(t(x))=(1+x)h(x)^2,
 \qquad t=\frac{x(1-8x)}{1+x},
 \label{eq:oracleB-franel-cover}
\end{equation}
and the quadratic extension with discriminant~$\Delta$.  It does not state
that $B_p$ is squarefree, and it does not formulate the theorem in terms of
a Hasse invariant, a Picard group, or an elliptic pencil.  The squarefree
and coprimality assertions in Remark~\ref{rem:squareness} therefore require
a separate proof or reference; the finite verification below is not such a
proof.

Second, the explicit pencil proposed in the oracle specification,
\begin{equation}
 E_t:\quad y^2=x(x-1)\bigl(x-t(1-t)\bigr),
 \label{eq:oracleB-proposed-pencil}
\end{equation}
is not identified with the Ap\'ery truncation in the cited material.  In
the displayed coordinate its discriminant is
\[
 16t^2(1-t)^2(1-t+t^2)^2,
\]
whose support differs from $\Delta(t)=0$.  More concretely, at $p=7$ its
raw Deuring polynomial is obtained by substituting
$\lambda=t(1-t)$ in
$1+2\lambda+2\lambda^2+\lambda^3$; its unique $\mathbb F_7$ root is
$t=4$, while the Ap\'ery polynomial above has root~$t=1$.  This does not
exclude an additional base change and gauge.  It proves that
\eqref{eq:oracleB-proposed-pencil}, as printed, cannot simply be declared
to produce~$B_p$, and in any event evaluation Frobenius cannot produce the
coefficient-first-zero~$\rho_p$ without the Mellin step
\eqref{eq:oracleB-mellin}.

\subsubsection{The moving-target quantifier trap}

Remark~\ref{rem:orbit} identifies $b_m\equiv0\pmod p$ with collisions in
the projective Ap\'ery orbit and supplies reflection.  For every fixed gap
$h$, the corresponding gap polynomial $N_h(m)\in\mathbb Z[m]$ is fixed,
so a prime-aspect Chebotarev theorem may be applied after~$h$ is fixed.
The first actual gap, however, may range up to $O(p)$, no uniform
$o(p)$ bound is known, and it varies with the same prime being counted.
The available quantifiers are
\[
 \text{for every fixed }h,\quad X\longrightarrow\infty;
\]
AMTD would require estimates uniform for moving $h\ll p$ inside
one dyadic block.  Degrees, splitting fields, discriminants, and effective
Chebotarev constants all move with~$h$.  Interchanging these limits is
exactly the forbidden moving-polynomial argument.

The center gives an even sharper warning.  Beukers' congruence in
Remark~\ref{rem:open} says
\[
 b_{(p-1)/2}\equiv a_p(f)\pmod p,
 \qquad f=\eta(2z)^4\eta(4z)^4.
\]
The index $(p-1)/2$ moves with~$p$, and the density-zero conjecture for
the nonordinary primes $p\mid a_p(f)$ is open for this non-CM weight-four
form.  Ordinary Sato--Tate controls fixed intervals for normalized traces;
it does not prove this shrinking, same-characteristic divisibility event.
It therefore cannot be invoked to prove B1 or B3.

Finally, equidistribution of $\rho_p$ modulo each fixed integer~$q$ would
still be only a marginal statement.  It has neither the uniformity in
growing frequencies needed for the minor arcs nor the two-prime joint
information needed for collision hyperplanes.  In particular it cannot,
by itself, exclude the scale-dependent anchored star.

\subsubsection{The correct fixed-family reduction}

The minimum horizontal statement can be written without referring to a
spurious root of~$B_p$.  For $N<p\leq2N$ and $m\in I_N=(N,2N]$, put
\[
 X_p(m)=\mathbf1_{\{p\leq m\}}
          \mathbf1_{\{b_{m-p}\equiv0\pmod p\}},
 \qquad
 T_N=\sum_{N<p\leq2N}\sum_{m\in I_N}X_p(m).
\]
Then the centered marked two-characteristic Hasse dispersion statement is
\begin{equation}
 \left|
 \sum_{\substack{N<p,q\leq2N\\p\ne q}}
 \sum_{m\in I_N}X_p(m)X_q(m)
 -\frac{T_N^2}{N}
 \right|
 \ll N^{o(1)}T_N.
 \label{eq:oracleB-mh2}
\end{equation}
Writing
\[
 F_N:=\sum_{m\in I_N}L_N(m)(L_N(m)-1),
\qquad L_N(m)=\sum_{N<p\leq2N}X_p(m),
\]
for the ordered collision count, the left side of
\eqref{eq:oracleB-mh2} is $|F_N-T_N^2/N|$.  The exact identity
\[
 \sum_{m\in I_N}\left(L_N(m)-\frac{T_N}{N}\right)^2
 =T_N+F_N-\frac{T_N^2}{N}
\]
shows directly that~\eqref{eq:oracleB-mh2} gives the top-block AMTD bound.
It also kills a full top-block anchored star (for example the anchor
$m_0=2N$), for which $L_N(m_0)$ has order $N/\log N$, and it is the
top-block form of the centered second-factorial-moment input used by AMTD.
The corresponding statement for all $P>\sqrt N$, with the centering of
Hypothesis~\ref{hyp:amtd}, is the precise input needed by the
analytic oracle.

A genuine fixed-family route would have to construct, from one integral
transcendental crystal (morally the symmetric square of the underlying
rank-two Picard--Fuchs connection), the Kummer-twisted objects
\begin{equation}
 \mathcal M_{p,j}=
 R\Gamma_{c,\mathrm{rig}}
 \bigl(U_p,\mathcal V_{\mathrm{tr},p}
       \otimes\mathcal K_{\omega_p^{-j}}\bigr)
 \label{eq:oracleB-crystalline-mellin}
\end{equation}
and a marked Hasse coordinate $c_{p,j}$ satisfying
\[
 c_{p,j}=0\quad\Longleftrightarrow\quad b_j\equiv0\pmod p.
\]
It would then have to prove~\eqref{eq:oracleB-mh2} for these marked zeros,
uniformly in the two characteristics and in the moving characters~$j$.
This is a $1/p$-scale crystalline anti-concentration and horizontal large
sieve statement, not ordinary Frobenius equidistribution.  Fixed
rank/conductor of an $\ell$-adic sheaf does not supply it: the Jacobi-sum
counterexample discussed with the Sym$^2$ remarks has fixed complexity but
linearly many bad tame characters.

Thus a construction of~\eqref{eq:oracleB-crystalline-mellin}, its marked
coordinate, and the centered two-characteristic estimate
\eqref{eq:oracleB-mh2} would be a quantifier-correct reduction of AMTD to a
fixed arithmetic family.  This is a specification of the missing package,
not an achieved G3 reduction to an existing Frobenius/Sato--Tate theorem:
neither the marked coordinate with all its integral control nor its
horizontal dispersion theorem is presently available.

\subsubsection{Finite verification}

The reproducible computation in \texttt{oracleB\_explore.py} recomputes
all $\mathcal Z_p$ for the $2260$ primes $5\leq p\leq20000$ and compares
them entrywise with the independent binary zero-set bank.  The exact
histogram is
\[
\begin{array}{c|rrrrrr}
 Z(p)&0&1&2&4&6&8\\ \hline
 \#p&1356&2&695&176&27&4.
\end{array}
\]
Among the 695 doublets, the KS distance of
$2\rho_p/(p-1)$ from the reflection-constrained uniform limit is
$0.03894$.  Small-modulus permutation tests for the outer gap
$h_p\bmod q$ against $p\bmod q$, using all $904$ active primes and
$q=3,5,7,11,13$, have p-values between $0.257$ and $0.675$.

An exact Sage/FLINT worker verified~\eqref{eq:oracleB-square}, the predicted
degree, squarefreeness, and coprimality to~$\Delta$ at every one of the
$2260$ primes.  It also verified $B_p\mid t^{p^2}-t$ and used
$\deg\gcd(B_p,t^p-t)$ to extract the exact linear/quadratic factorization
type at every prime in the range.  Complete factorization into explicit
factor polynomials was restricted to a sparse cross-check sample.
At $p=101,503,1009,5003,10007,19997$, the degrees of $B_p$ are respectively
$50,250,504,2501,5002,9998$, and the numbers of irreducible factors are
$30,148,262,1284,2595,5058$.  These calculations demonstrate growing
complexity and catch transcription errors; they prove no assertion beyond
the finite range.  Exact identity checks used in the propositions above are
repeated independently by \texttt{oracleB\_verify.py}.

\paragraph{STALL REPORT.}
The actual cross-prime decorrelation of the Ap\'ery coefficient zero-sets,
and hence G1/AMTD, is not proved.  G3 is not reached for the reason stated
above; literal G2 is proved but has insufficient uniformity; the obstruction
and exact missing input constitute G4.  The rigorous outputs are:
\begin{enumerate}
 \item Proposition~\ref{prop:oracleB-fixed-anchor}, which proves the
       literal G2 count is $O_{m_0}(1)$ for every fixed anchor;
 \item Remark~\ref{rem:oracleB-quantifiers}, which proves that this has
       the wrong quantifier order for the moving anchored-star family;
 \item Proposition~\ref{prop:oracleB-two-zero-loci} and the exact $p=7$
       counterexample, which show that B1/B2 as stated confuse Mellin
       coefficient zeros with evaluation roots;
 \item Proposition~\ref{prop:oracleB-complexity}, which isolates the
       linear-degree obstruction even on the evaluation-root side;
 \item the fixed-gap/fixed-prime-family audit, including the open
       nonordinary center slice from Remark~\ref{rem:open};
 \item the quantifier-correct remaining arithmetic input
       \eqref{eq:oracleB-mh2}, together with the crystalline Mellin package
       \eqref{eq:oracleB-crystalline-mellin} that a fixed-family proof would
       have to supply.
\end{enumerate}
